\chapter{Prawda arytmetyczna i twierdzenie Tarskiego}\label{ch:tarski-truth}

\section{Prawda to nie to samo co dowodliwość}
W rozdziale~\ref{ch:first-incompleteness} udowodniliśmy,
że dostatecznie silna, niesprzeczna, efektywnie
aksjomatyzowalna teoria arytmetyczna nie rozstrzyga
wszystkich zdań. Możemy jednak zapytać o jeszcze
silniejszy pomysł: czy istnieje \emph{jedna formuła
arytmetyczna}, która dla każdej liczby kodującej
zdanie rozpoznaje, czy to zdanie jest prawdziwe
w standardowej strukturze liczb naturalnych?

Odpowiedź jest przecząca. Zanim ją udowodnimy,
oddzielmy trzy poziomy języka:
\begin{itemize}
\item Zdanie arytmetyczne $\sigma$ jest \emph{napisem}.
\item Liczba $s=\ulcorner\sigma\urcorner$ jest jego kodem,
      a $\overline s$ termem będącym numerałem kodu.
\item Zapis $\mathfrak N\models\sigma$ jest
      \emph{metamatematycznym} stwierdzeniem prawdziwości
      tego zdania w zwykłych liczbach naturalnych.
\end{itemize}
Sama notacja $\mathfrak N\models\sigma$ nie jest dodatkowym
predykatem języka $\{0,S,+,\cdot,<\}$.

\begin{przyklad}[Dwie wypowiedzi o liczbie $4$]
Zdanie $\overline2+\overline2=\overline4$ jest prawdziwe
w $\mathfrak N$, czyli
$\mathfrak N\models(\overline2+\overline2=\overline4)$.
Niech $T\supseteq Q$ będzie efektywnie aksjomatyzowalną
teorią i niech $g$ oznacza kod tego zdania. Inne zdanie,
$\operatorname{Bew}_T(\overline g)$, mówi o istnieniu
\emph{dowodu} danej równości w $T$.
Ponieważ $Q$ dowodzi tej równości, istnieje jej rzeczywisty,
skończony dowód w $T$. Zatem oba zdania są prawdziwe
w $\mathfrak N$, niezależnie od tego, czy $T$ jest poprawna
semantycznie. Poprzedni rozdział pokazał jednak, że taka
zgodność nie zachodzi w ogólności dla wszystkich zdań
arytmetycznych.
\end{przyklad}

\section{Co miałoby znaczyć „formuła definiująca prawdę”?}
\begin{definicja}[Definiowalność prawdy arytmetycznej]
Formuła $\operatorname{Tr}(x)$ języka arytmetyki,
w której jedyną zmienną mogącą występować wolno jest $x$,
\emph{definiowałaby prawdę arytmetyczną w $\mathfrak N$},
gdyby dla każdego \emph{zdania} $\sigma$ tego języka
zachodziła równoważność metamatematyczna
\[
\mathfrak N\models\operatorname{Tr}(\overline{\ulcorner\sigma\urcorner})
\quad\Longleftrightarrow\quad
\mathfrak N\models\sigma.
\]
Nie narzucamy żadnego wymagania na wartości
$\operatorname{Tr}(x)$ dla liczb niekodujących zdań.
Równoważność dotyczy więc wyłącznie kodów \emph{zdań},
a nie dowolnych napisów lub formuł z wolnymi zmiennymi.
\end{definicja}

\begin{uwaga}[Dwa różne rodzaje równoważności]
Jeśli $s=\ulcorner\sigma\urcorner$, zdanie
$\sigma\leftrightarrow\operatorname{Tr}(\overline s)$
można \emph{zapisać} w języku arytmetyki.
Natomiast stwierdzenie
$\mathfrak N\models(\sigma\leftrightarrow\operatorname{Tr}(\overline s))$
jest wypowiedzią metamatematyczną o prawdziwości tej równoważności.
Teoria $T$ może również \emph{dowodzić}
$T\vdash\sigma\leftrightarrow\operatorname{Tr}(\overline s)$.
Są to dwa różne wymagania. Dopiero gdy wszystkie twierdzenia $T$
są prawdziwe w $\mathfrak N$ (czyli gdy $T$ jest poprawna
względem $\mathfrak N$), dowodliwość pociąga za sobą
prawdziwość w $\mathfrak N$. W ogólności takiego przejścia
nie wolno wykonywać.
\end{uwaga}

\section{Twierdzenie Tarskiego}
\begin{twierdzenie}[Niedefiniowalność prawdy --- wersja semantyczna]\label{thm:tarski-semantic}
Nie istnieje formuła $\operatorname{Tr}(x)$ języka
$\mathcal L_{\mathrm{ar}}=\{0,S,+,\cdot,<\}$,
która definiuje prawdę wszystkich zdań tego języka
w standardowej strukturze $\mathfrak N$.
\end{twierdzenie}
\begin{proof}
Przypuśćmy przeciwnie, że taka formuła istnieje.
Stosujemy lemat diagonalny
(Lemat~\ref{res:8:1}) do formuły
$\psi(x):=\neg\operatorname{Tr}(x)$.
Dostajemy zdanie arytmetyczne $L$ takie, że
\[
Q\vdash L\leftrightarrow\neg\operatorname{Tr}(\overline\ell),
\qquad \ell:=\ulcorner L\urcorner.
\]
Ponieważ $\mathfrak N$ jest modelem $Q$, w modelu
standardowym zachodzi
\begin{equation}\label{eq:tarski-anti}
\mathfrak N\models
L\leftrightarrow\neg\operatorname{Tr}(\overline\ell).
\end{equation}
Z założenia, że $\operatorname{Tr}$ definiuje prawdę,
dla \emph{tego samego} zdania $L$ dostajemy również
\begin{equation}\label{eq:tarski-truth}
\mathfrak N\models
L\leftrightarrow\operatorname{Tr}(\overline\ell).
\end{equation}

W klasycznej logice prawdziwość~\eqref{eq:tarski-truth}
oznacza, że $L$ i $\operatorname{Tr}(\overline\ell)$
mają tę samą wartość logiczną. Tymczasem
\eqref{eq:tarski-anti} mówi, że wartości te są przeciwne.
Jeśli $L$ jest prawdziwe, to z~\eqref{eq:tarski-anti}
wynika fałszywość $\operatorname{Tr}(\overline\ell)$,
a z~\eqref{eq:tarski-truth} jej prawdziwość.
Jeśli $L$ jest fałszywe, z pierwszej równoważności
wynika prawdziwość $\operatorname{Tr}(\overline\ell)$,
a z drugiej jej fałszywość. Oba przypadki są niemożliwe.
Zatem formuła $\operatorname{Tr}$ nie istnieje.
\end{proof}

\begin{uwaga}[Dlaczego to nie jest zwykły paradoks kłamcy?]
Nie wpisaliśmy do języka arytmetycznego nieformalnego
zdania „to zdanie jest fałszywe”.
Najpierw kodujemy skończone napisy, następnie
konstruujemy funkcję podstawienia na kodach,
a na koniec używamy \emph{udowodnionego}
Lematu~\ref{res:8:1}. Zdanie $L$ i jego kod $\ell$
to różne obiekty. Sprzeczność powstaje dopiero
po dołożeniu hipotezy, że jedna wewnętrzna formuła
rozpoznaje prawdę \emph{wszystkich} zdań.
\end{uwaga}

\begin{twierdzenie}[Wersja syntaktyczna Tarskiego]\label{thm:tarski-syntactic}
Niech $T$ będzie dowolną niesprzeczną teorią
klasycznej logiki pierwszego rzędu zawierającą $Q$.
Nie istnieje formuła arytmetyczna $\operatorname{Tr}(x)$
taka, że dla \emph{każdego} zdania arytmetycznego
$\sigma$ teoria $T$ dowodzi
\[
\sigma\leftrightarrow
\operatorname{Tr}(\overline{\ulcorner\sigma\urcorner}).
\]
Nie zakładamy efektywnego wyliczania aksjomatów $T$.
\end{twierdzenie}
\begin{proof}
Załóżmy, że wszystkie takie równoważności są dowodliwe.
Dla $\psi(x)=\neg\operatorname{Tr}(x)$ lemat diagonalny
daje zdanie $L$ i liczbę $\ell=\ulcorner L\urcorner$ z
\[
T\vdash L\leftrightarrow\neg\operatorname{Tr}(\overline\ell),
\]
ponieważ $T\supseteq Q$. Założenie zastosowane
\emph{do zdania $L$} daje jednocześnie
$T\vdash L\leftrightarrow\operatorname{Tr}(\overline\ell)$.
Złożenie obu równoważności daje
$T\vdash L\leftrightarrow\neg L$.
Ten krok jest czysto syntaktyczny: nie zakłada, że zdania
dowodliwe w $T$ są prawdziwe w $\mathfrak N$.
Wyprowadźmy jawnie sprzeczność.
Z kierunku $L\to\neg L$ i klasycznego rozumowania
przez sprzeczność otrzymujemy $T\vdash\neg L$.
Z drugiego kierunku $\neg L\to L$ otrzymujemy następnie
$T\vdash L$. Zatem $T$ dowodzi zdania i jego negacji,
co przeczy niesprzeczności.
\end{proof}

\begin{przyklad}[Dlaczego problem dotyczy jednego uniwersalnego predykatu?]
Dla ustalonej skończonej, niepustej listy zdań
$\sigma_1,\ldots,\sigma_n$ (gdzie $n\ge1$)
możemy ustalić kody $g_i=\ulcorner\sigma_i\urcorner$ i zapisać
\[
\operatorname{Tr}_{\mathrm{fin}}(x):=
\bigvee_{i=1}^{n}(x=\overline{g_i}\land\sigma_i).
\]
Dla każdego $i\in\{1,\ldots,n\}$ w $\mathfrak N$ zachodzi
$\operatorname{Tr}_{\mathrm{fin}}(\overline{g_i})
\leftrightarrow\sigma_i$.
Jeżeli na liście powtórzono to samo zdanie, jego powtórzenia
mają identyczną wartość logiczną. Natomiast różne zdania
mają różne kody dzięki różnowartościowości numeracji;
dla danego $g_i$ wszystkie człony o innych kodach są fałszywe.
Nie otrzymaliśmy jednak \emph{jednej} formuły
obsługującej równocześnie wszystkie zdania arytmetyczne:
po dodaniu nowych zdań do listy musimy zmienić formułę.
Twierdzenie Tarskiego nie przeczy takiej skończonej konstrukcji.
\end{przyklad}

\section{Dlaczego prawda nie jest tym samym co dowodliwość?}
\begin{twierdzenie}[Niesprzeczna teoria nie dowodzi równoważności prawdy i dowodliwości]
Niech $T\supseteq Q$ będzie niesprzeczna i efektywnie
aksjomatyzowalna. Wówczas nie jest prawdą, że dla każdego
zdania $\sigma$ zachodzi
$T\vdash\sigma\leftrightarrow
\operatorname{Bew}_T(\overline{\ulcorner\sigma\urcorner})$.
\end{twierdzenie}
\begin{proof}
Gdyby wszystkie te równoważności były dowodliwe,
formuła $\operatorname{Tr}(x):=\operatorname{Bew}_T(x)$
spełniałaby warunek wykluczony w
Twierdzeniu~\ref{thm:tarski-syntactic}. Ponieważ $T$
jest niesprzeczna i zawiera $Q$, otrzymujemy sprzeczność.
\end{proof}

\begin{uwaga}[Dwa rodzaje ograniczeń]
Pierwsze twierdzenie Gödla mówi o granicach
\emph{dowodzenia} w pojedynczej efektywnie aksjomatyzowalnej,
niesprzecznej teorii rozszerzającej $Q$.
Semantyczne twierdzenie Tarskiego mówi o niemożliwości
\emph{zdefiniowania pełnej prawdy} arytmetycznej
jedną formułą tego samego języka i nie wymaga wyboru
teorii $T$. W jego wersji syntaktycznej można wybrać
dowolną niesprzeczną teorię $T\supseteq Q$,
bez założenia efektywnej aksjomatyzowalności. Nie są to identyczne twierdzenia,
choć oba korzystają z tego samego mechanizmu diagonalizacji.
\end{uwaga}

\section{Zapowiedź warunków wyprowadzalności}
Do badania drugiego twierdzenia o niezupełności
nie wystarczy znajomość tego, które konkretne
instancje predykatu $\operatorname{Prf}_T$
są prawdziwe w $\mathfrak N$.
Potrzebujemy formalnych relacji między zdaniami
mówiącymi o \emph{istnieniu dowodów}.

\begin{definicja}[Skrót $\Box\varphi$]\label{def:box-provability}
Ustalmy efektywnie aksjomatyzowalną teorię
$T\supseteq\mathrm{PA}$ i jej predykat
$\operatorname{Bew}_T(x)$.
Dla \emph{zdania} arytmetycznego $\varphi$ oznaczmy
$g=\ulcorner\varphi\urcorner$.
Zapis $\Box\varphi$ jest skrótem pełnego zdania
$\operatorname{Bew}_T(\overline g)$.
Jeśli chcemy zapisać $\Box\Box\varphi$, najpierw
tworzymy zdanie $\Box\varphi$, potem bierzemy
\emph{jego} kod $h=\ulcorner\Box\varphi\urcorner$
i zapisujemy $\operatorname{Bew}_T(\overline h)$.
Nie podnosimy liczby $g$ do kwadratu: drugi znak $\Box$
oznacza nowy poziom wypowiedzi o dowodzie.
Jest to skrót \emph{metamatematyczny}, a nie nowy
spójnik dopisany do alfabetu arytmetyki.
\end{definicja}

\begin{definicja}[Warunki Hilberta--Bernaysa--Löba]\label{def:hilbert-bernays-loeb}
Niech $T\supseteq\mathrm{PA}$ będzie efektywnie
aksjomatyzowalną teorią klasycznej logiki pierwszego rzędu,
a $\Box$ ustalonym predykatem dowodliwości,
rozumianym jak w poprzedniej definicji.
Mówimy, że spełnia on następujące warunki, jeśli dla
wszystkich zdań arytmetycznych $\varphi,\psi$ zachodzą:
\begin{align*}
\textup{(D1)}\quad&T\vdash\varphi\ \Longrightarrow\ T\vdash\Box\varphi,\\
\textup{(D2)}\quad&T\vdash\Box(\varphi\to\psi)\to(\Box\varphi\to\Box\psi),\\
\textup{(D3)}\quad&T\vdash\Box\varphi\to\Box\Box\varphi.
\end{align*}
Pierwsza linia jest regułą o \emph{rzeczywistych dowodach};
pozostałe dwie linie są \emph{formułami dowodliwymi w $T$}.
\end{definicja}

\begin{lemat}[Warunek (D1) dla standardowego predykatu]\label{lem:d1-provability}
Niech $T\supseteq\mathrm{PA}$ będzie efektywnie
aksjomatyzowalna, a $\Box$ pochodzi z predykatu
$\operatorname{Prf}_T$ z Lematu~\ref{res:7:7}.
Wtedy zachodzi warunek \textup{(D1)}.
\end{lemat}
\begin{proof}
Jeżeli $T\vdash\varphi$, wybieramy jeden
rzeczywisty kod $p$ dowodu $\varphi$.
Niech $g=\ulcorner\varphi\urcorner$.
Ponieważ $p$ jest rzeczywistym kodem poprawnego dowodu,
zdanie $\operatorname{Prf}_T(\overline p,\overline g)$
jest prawdziwą zamkniętą instancją formuły $\Delta_0$.
Z części (i) Lematu~\ref{lem:incomp-proof-check}
wynika więc $Q\vdash\operatorname{Prf}_T(\overline p,\overline g)$.
Wprowadzając kwantyfikator egzystencjalny,
dostajemy $Q\vdash\Box\varphi$.
Ponieważ $T\supseteq\mathrm{PA}\supseteq Q$,
ten sam wniosek jest dowodliwy w $T$.
\end{proof}

\begin{przyklad}[Dlaczego (D1) nie działa w drugą stronę?]
Warunek (D1) mówi: jeśli \emph{mamy konkretny dowód}
$\varphi$, to możemy w teorii udowodnić,
że jakiś dowód $\varphi$ istnieje. Nie mówi,
że z $T\vdash\Box\varphi$ zawsze wolno wyciągnąć
wniosek $T\vdash\varphi$.
Z samego dowodu zdania egzystencjalnego
$\exists p\,\operatorname{Prf}_T(p,\overline g)$
nie wynika możliwość wskazania \emph{standardowej liczby
naturalnej} kodującej dowód $\varphi$: w niestandardowym
modelu arytmetyki świadkiem tej formuły może być element,
który nie jest kodem rzeczywistego, skończonego dowodu.
Dodatkowe założenie może zmienić sytuację. Jeśli $T$ jest
$\Sigma_1$-poprawna względem $\mathfrak N$ (to znaczy:
każde zdanie $\Sigma_1$ dowodliwe w $T$ jest prawdziwe
w $\mathfrak N$), to z
$T\vdash\Box\varphi$ wynika
$\mathfrak N\models\Box\varphi$, gdyż
$\Box\varphi$ jest zdaniem $\Sigma_1$.
Wtedy istnieje standardowa liczba $p$, dla której
$\mathfrak N\models\operatorname{Prf}_T(\overline p,\overline g)$,
a więc rzeczywisty dowód $\varphi$ w $T$.
Takiej $\Sigma_1$-poprawności nie zakłada jednak (D1);
jego odwrócenie nie wynika z samego warunku wyprowadzalności.
\end{przyklad}

\begin{uwaga}[Co pozostaje do udowodnienia?]
Reguły \textup{(D2)} i \textup{(D3)} wymagają
\emph{wewnętrznego} kodowania operacji łączenia dowodów
i weryfikacji kodu dowodu.
Nie wynika z (D1), że są automatycznie spełnione.
W rozdziale~\ref{ch:derivability-proof} udowodnimy je
dla odpowiedniego standardowego predykatu w arytmetyce
Peana, zamiast przyjmować je bez uzasadnienia.
\end{uwaga}

\section{Jak czytać trzy warunki na konkretnym przykładzie?}
\begin{przyklad}[Od rachunku $1+1=2$ do zdania o dowodzie]
Niech $\varphi$ będzie zdaniem
$S(0)+S(0)=S(S(0))$ i niech $T\supseteq\mathrm{PA}$
będzie efektywnie aksjomatyzowalna. Z aksjomatów
dodawania otrzymujemy kolejno
\[
S(0)+S(0)=S(S(0)+0)=S(S(0)),
\]
a więc $T\vdash\varphi$. Możemy wybrać kod $p$
\emph{jednego konkretnego} dowodu tej równości.
Weryfikacja $p$ daje
$Q\vdash\operatorname{Prf}_T(\overline p,\overline g)$,
gdzie $g=\ulcorner\varphi\urcorner$.
Po wprowadzeniu kwantyfikatora egzystencjalnego
uzyskujemy $T\vdash\Box\varphi$.
Tak działa \textup{(D1)}. Nie używaliśmy tu wniosku,
że z istnienia kodu dowodu wynika prawdziwość dowodzonego
zdania w każdym możliwym modelu; posłużyliśmy się
\emph{rzeczywistym dowodem}, który już mieliśmy.
\end{przyklad}

\begin{uwaga}[Dlaczego (D2) jest formalizacją modus ponens?]
Reguła \emph{modus ponens} pozwala z dowodów
$\varphi\to\psi$ oraz $\varphi$ otrzymać dowód $\psi$.
Na poziomie kodów odpowiada jej obliczalna operacja,
która skleja dwa skończone dowody i dopisuje
ostatni wiersz. Warunek \textup{(D2)} mówi,
że teoria $T$ potrafi uzasadnić istnienie
odpowiedniego kodu \emph{wewnątrz arytmetyki}:
\[
\Box(\varphi\to\psi)\to(\Box\varphi\to\Box\psi).
\]
Zwróćmy uwagę na nawiasy: jest to \emph{jedno zdanie}
$T$ i nie należy go mylić z regułą o \emph{dwóch
gotowych dowodach}.
\end{uwaga}

\begin{uwaga}[Dlaczego (D3) nie znaczy, że teoria jest nieomylna?]
Zdanie $\Box\varphi$ głosi, że istnieje kod
dowodu $\varphi$. Warunek \textup{(D3)} głosi,
że z takiego istnienia w teorii wynika istnienie
kodu dowodu \emph{tego właśnie zdania o dowodzie}:
\[
\Box\varphi\to\Box\Box\varphi.
\]
Nie jest to implikacja $\Box\varphi\to\varphi$.
Ta ostatnia byłaby \emph{wewnętrznym stwierdzeniem
poprawności dowodów} $\varphi$, a nie jedynie zdaniem
o kolejnych dowodach. Nie wolno jej utożsamiać
z metamatematycznym warunkiem prawdziwości $\varphi$
w $\mathfrak N$.
Warunek (D3) dotyczy wyłącznie dalszego
poziomu \emph{dowodliwości}. W następnym rozdziale
przedstawimy formalne uzasadnienie (D2) i (D3);
sama definicja warunków nie gwarantuje ich spełnienia.
\end{uwaga}


\section{Zakres twierdzenia Tarskiego}
Twierdzenie Tarskiego dotyczy \emph{jednej} formuły arytmetycznej,
która miałaby poprawnie rozpoznawać wszystkie zdania pierwszego
rzędu w standardowym modelu. Nie wyklucza definiowania prawdy
dla formuł o ustalonej, ograniczonej złożoności ani zewnętrznej
definicji spełniania w bogatszym metajęzyku. Również predykat
$\operatorname{Bew}_T$ nie jest predykatem prawdy: koduje istnienie
formalnego dowodu, którego poprawność jest syntaktyczna.

\section{Definiowalność, wyliczalność i rozstrzygalność to różne własności}
\begin{uwaga}
\emph{Definiowalność} podzbioru $A\subseteq\N$
oznacza istnienie formuły arytmetycznej $\theta(x)$
bez innych zmiennych wolnych, takiej że dla każdego $n\in\N$
zachodzi $n\in A$ wtedy i tylko wtedy, gdy
$\mathfrak N\models\theta(\overline n)$. \emph{Efektywna wyliczalność}
oznacza istnienie programu wypisującego wszystkie jego elementy i tylko je.
\emph{Rozstrzygalność} oznacza istnienie programu, który dla każdego
argumentu kończy działanie z poprawną odpowiedzią „tak” albo „nie”.
Nie są to pojęcia równoważne. Zbiór twierdzeń efektywnie aksjomatyzowalnej
teorii można efektywnie wyliczać, lecz na mocy wyników z rozdziału
\ref{ch:first-incompleteness} dla efektywnie aksjomatyzowalnych, niesprzecznych
rozszerzeń $Q$ nie można go rozstrzygać.
Pełna prawda arytmetyczna nie jest nawet definiowalna jedną formułą
tego samego języka, co wynika z twierdzenia Tarskiego.

Dla zbioru nieskończonego przeliczalność oznacza istnienie
bijekcji z $\N$, natomiast określenie „co najwyżej przeliczalny”
obejmuje również zbiory skończone i odpowiada istnieniu
iniekcji do $\N$. Żaden z tych warunków sam w sobie
nie dostarcza algorytmu wyznaczającego taką bijekcję lub iniekcję. W kolejnych rozdziałach
\ref{ch:downward-ls} i \ref{ch:upward-ls} zobaczymy,
że konstruowanie modeli przeliczalnych może wymagać wyborów
świadków, których prawdziwości nie umiemy algorytmicznie sprawdzić.
\end{uwaga}

\begin{przyklad}[Dlaczego arytmetyzacja nie tworzy predykatu prawdy?]
Kodujemy napis zdania jako liczbę naturalną $n$,
a więc możemy pytać o jego składnię i o istnienie formalnego
dowodu tego zdania. Nie otrzymujemy jednak automatycznie
algorytmu ani formuły mówiącej, czy zakodowane zdanie
jest prawdziwe w standardowej arytmetyce.
Należy odróżniać badanie skończonego kodu składni
od prawdziwości zdania w standardowej strukturze liczb naturalnych.
Dlatego diagonalizacja nie czyni arytmetyki sprzeczną:
zdanie Gödla mówi o dowodliwości, natomiast
niedozwolony uniwersalny predykat Tarskiego miałby
wyrażać prawdę własnego języka.
\end{przyklad}